\chapter{✨ 形容詞 — Adjetivos}
\unidadheader{09}{形容詞}{Adjetivos}{Describir personas, lugares y cosas}

\begin{tcolorbox}[vocabbox, title={📌 Vocabulario Nuevo}]
\begin{tabularx}{\linewidth}{llX}
\toprule
\textbf{Japonés} & \textbf{Español} & \textbf{Tipo} \\
\midrule
\vocabitem{きれい}{—}{bonito/a / limpio/a}{adj-な}
\vocabitem{面白い}{おもしろい}{interesante / gracioso}{adj-い}
\vocabitem{つまらない}{—}{aburrido}{adj-い}
\vocabitem{便利}{べんり}{conveniente / práctico}{adj-な}
\vocabitem{新しい}{あたらしい}{nuevo/a}{adj-い}
\vocabitem{古い}{ふるい}{viejo/a / antiguo}{adj-い}
\vocabitem{難しい}{むずかしい}{difícil}{adj-い}
\vocabitem{やさしい}{—}{fácil / amable}{adj-い}
\vocabitem{楽しい}{たのしい}{divertido / entretenido}{adj-い}
\vocabitem{暇}{ひま}{libre (sin plan) / aburrimiento}{adj-な/sust.}
\bottomrule
\end{tabularx}
\end{tcolorbox}

\begin{tcolorbox}[phrasebox, title={⚙️ Conjugación de adjetivos}]
\begin{tabular}{lllll}
\toprule
\textbf{Tipo} & \textbf{Presente +} & \textbf{Presente -} & \textbf{Pasado +} & \textbf{Pasado -} \\
\midrule
adj-い & 〜い です & 〜くない & 〜かった & 〜くなかった \\
adj-な & 〜な です & 〜じゃない & 〜でした & 〜じゃなかった \\
\bottomrule
\end{tabular}\\[4pt]
\small Ej: \ruby{面白}{おもしろ}い → \ruby{面白}{おもしろ}くない / \ruby{面白}{おもしろ}かった / \ruby{面白}{おもしろ}くなかった
\end{tcolorbox}

\subsection*{🗣️ Frases Clave}

\frase{この\ruby{映画}{えいが}は\ruby{面白}{おもしろ}いですか？}
  {¿Es interesante esta película?}
\frase{あまり\ruby{面白}{おもしろ}くなかったです。}
  {No fue muy interesante.}
\frase{この\ruby{町}{まち}はとても\ruby{便利}{べんり}です。}
  {Esta ciudad es muy conveniente.}
\frase{\ruby{日本語}{にほんご}は\ruby{難}{むずか}しいですか？}
  {¿Es difícil el japonés?}

\subsection*{⚙️ Gramática}

\patron{Adjetivos い — todas las formas}
  {adj-い: 〜い / 〜くない / 〜かった / 〜くなかった}
  {\ej{この\ruby{本}{ほん}は\ruby{面白}{おもしろ}いです。}{Este libro es interesante.}
  \ej{この\ruby{本}{ほん}は\ruby{面白}{おもしろ}くないです。}{Este libro no es interesante.}
  \ej{きのうの\ruby{授業}{じゅぎょう}は\ruby{難}{むずか}しかったです。}
    {La clase de ayer fue difícil.}}

\patron{Adjetivos な — antes de sustantivo}
  {adj-な + な + 名詞}
  {\ej{\ruby{便利}{べんり}なアプリです。}{Es una aplicación práctica.}
  \ej{きれいな\ruby{花}{はな}ですね。}{Es una flor bonita, ¿verdad?}}

\begin{tcolorbox}[ejerciciobox, title={📝 Ejercicios Rápidos}]
\begin{enumerate}
  \item Conjuga en pasado negativo: \ruby{楽}{たの}しい → ?
  \item ¿Cuál es la diferencia entre やさしい y きれい como tipo de adjetivo?
  \item Describe tu ciudad usando dos adjetivos.
\end{enumerate}
\vspace{0.4em}
\textbf{Respuestas:}
1) \ruby{楽}{たの}しくなかった(です)\quad
2) やさしい = adj-い; きれい = adj-な (no termina en い fonético para conjugar)\quad
3) (personal)
\end{tcolorbox}

\begin{tcolorbox}[kanjibox, title={🀄 Kanjis de la Unidad}]
\begin{tabularx}{\linewidth}{cllXX}
\toprule
\textbf{Kanji} & \textbf{On} & \textbf{Kun} & \textbf{Significado} & \textbf{Vocab} \\
\midrule
\kanjirow{新}{しん}{あたら・あら}{nuevo}{\ruby{新しい}{あたらしい} / \ruby{新幹線}{しんかんせん}}
\kanjirow{古}{こ}{ふる}{antiguo / viejo}{\ruby{古い}{ふるい} / \ruby{古典}{こてん}}
\kanjirow{難}{なん・なん}{むずか}{difícil}{\ruby{難しい}{むずかしい} / \ruby{難問}{なんもん}}
\kanjirow{便}{べん}{たよ}{conveniente}{\ruby{便利}{べんり} / \ruby{便}{べん}}
\kanjirow{楽}{がく・らく}{たの}{placer / fácil}{\ruby{楽しい}{たのしい} / \ruby{楽}{らく}}
\bottomrule
\end{tabularx}
\end{tcolorbox}
